\cleardoublepage

\chapter{KI in der Wissenschaft – Werkzeug statt Wahrheit}
\label{anhangC:ki}
\markboth{Anhang C}{Anhang C}

\subsection*{Motivation}

Dieses Buch entstand aus dem Wunsch, einen Gedanken konsequent auszuarbeiten, der sich durch die gesamte Geschichte der Mathematik und Physik zieht: Die Welt ist nicht bloß vorhanden, sondern strukturiert -- und Mathematik ist die präziseste Form, diese Struktur zu erkennen. Ziel des Buches ist es daher, Mathematik nicht nur als Sammlung von Techniken, sondern als innere Ordnung von Zahl, Form, Bewegung, Zufall und Natur verständlich zu machen.

Dabei wurde ein neues Werkzeug eingesetzt, das heute immer stärker in wissenschaftliches und schriftstellerisches Arbeiten hineinwirkt: \textbf{künstliche Intelligenz}\index{Künstliche Intelligenz}, konkret das Sprachmodell \textbf{ChatGPT}\index{ChatGPT} von OpenAI.

Doch wie lässt sich KI sinnvoll in einem Buchprojekt nutzen, das gerade auf begriffliche Klarheit, mathematische Genauigkeit und geistige Eigenständigkeit angewiesen ist? Und wie kann man ihren Einsatz offenlegen, ohne die Verantwortung des Autors zu verwischen? Dieser Anhang gibt einen transparenten Einblick in den Entstehungsprozess dieses Buches und plädiert für einen verantwortungsvollen Umgang mit KI als Werkzeug -- nicht als Wahrheit.

\subsection*{Was eine KI kann -- und was nicht}

KI-gestützte Sprachmodelle wie ChatGPT sind leistungsfähige Hilfsmittel beim Schreiben und Strukturieren. Sie können:

\begin{itemize}
	\item beim Formulieren erster Entwürfe helfen,
	\item komplexe Sachverhalte sprachlich glätten,
	\item Denkanstöße liefern oder Gliederungen vorschlagen,
	\item stilistische Alternativen aufzeigen.
\end{itemize}

Was sie jedoch \textbf{nicht} können:

\begin{itemize}
	\item \textbf{verstehen}\index{Verstehen} im wissenschaftlichen Sinn,
	\item \textbf{prüfen}, ob eine Formel korrekt hergeleitet ist,
	\item \textbf{physikalische Konzepte durchdringen},
	\item \textbf{Quellen kritisch einordnen oder bewerten}.
\end{itemize}

Daher gilt: Eine KI kann \emph{unterstützen}, aber sie \textbf{kann und darf den wissenschaftlichen Erkenntnisprozess nicht ersetzen}. Wer mit KI arbeitet, muss dennoch selbst denken -- und jedes Ergebnis kritisch prüfen.

\subsection*{Wie dieses Buch entstanden ist}

Die Inhalte dieses Buches -- von der Struktur über die physikalischen Erklärungen bis zu den mathematischen Herleitungen -- wurden vom Autor konzipiert, recherchiert und verantwortet. ChatGPT kam in folgenden Bereichen unterstützend zum Einsatz:

\begin{itemize}
	\item beim \textbf{Formulieren einzelner Passagen}, etwa bei Einleitungen, Zusammenfassungen oder didaktischen Abschnitten,
	\item zur \textbf{stilistischen Überprüfung} technischer Abschnitte,
	\item zur \textbf{Entwicklung von Gliederungen} in frühen Arbeitsphasen,
	\item zur Reflexion über \textbf{Verständlichkeit}\index{Verständlichkeit} und Leserführung.
\end{itemize}

Entscheidend ist: \textbf{Alle inhaltlichen Aussagen, Formeln und Interpretationen wurden vom Autor geprüft, hinterfragt, überarbeitet oder verworfen.} Keine KI ersetzte die fachliche Entwicklung der physikalischen Argumentation. Die inhaltliche Verantwortung, Auswahl und Bewertung aller Aussagen lag durchgehend beim Autor.

\subsection*{Ethische Fragen und wissenschaftliche Verantwortung}

Die Nutzung von KI in der Wissenschaft wirft berechtigte Fragen auf:

\begin{itemize}
	\item Wie viel darf automatisiert entstehen, ohne dass Autorschaft\index{Autorschaft} verwässert?
	\item Wie geht man mit potenziellen Fehlern\index{Fehler} um?
	\item Wie transparent muss die Nutzung offengelegt werden?
\end{itemize}

Die Antwort liegt in einem Grundprinzip wissenschaftlicher Arbeit: \textbf{Verantwortung}. Wer KI einsetzt, bleibt für das Ergebnis verantwortlich -- unabhängig davon, ob einzelne Formulierungen von einem Modell vorgeschlagen wurden.

In diesem Sinn ist KI keine Autorin, sondern ein Werkzeug. Sie kann Prozesse beschleunigen, aber nicht ersetzen, was Wissenschaft im Kern ausmacht: \textbf{kritisches Denken, sorgfältiges Prüfen und methodisches Arbeiten}.

\subsection*{Empfehlungen für den Einsatz in der Forschung}

Für Forschende, Lehrende und Studierende ergibt sich daraus ein konstruktiver Weg:

\begin{itemize}
	\item Nutze KI \textbf{bewusst und gezielt} -- für sprachliche Unterstützung, nicht für Argumentation oder Beweisführung.
	\item \textbf{Prüfe jede Aussage selbst} -- gerade bei komplexen Sachverhalten.
	\item \textbf{Deklariere die Nutzung offen}, wenn es relevant ist -- etwa in Vorworten, Anhängen oder Einreichungserklärungen.
	\item Nutze KI nicht zur \textbf{Täuschung}\index{Täuschung} oder als Feigenblatt, sondern als Hilfe zur besseren Darstellung eigener Gedanken.
\end{itemize}

\subsection*{Fazit: KI als Werkzeug -- aber der Mensch bleibt denkend verantwortlich}

Künstliche Intelligenz ist weder Ersatz noch Gegner menschlicher Erkenntnis. Sie ist ein \textbf{Werkzeug}\index{Werkzeug}, das bei der wissenschaftlichen Kommunikation helfen kann -- \textbf{wenn es bewusst, reflektiert und verantwortungsvoll eingesetzt wird}.

Dieses Buch versteht sich auch in dieser Hinsicht als Beitrag zu einem neuen, aufgeklärten Umgang mit Technologie in der Wissenschaft. Nicht, weil die Technik alles kann, sondern weil wir lernen müssen, sie sinnvoll zu nutzen.

\vspace{1em}

\section*{Hinweis zur Open-Access-Version}

Dieses Buch ist Teil der Open-Science-Initiative „Christian \& Co-Pilot -- Math \& Physics“.

Die vollständige, farbige PDF-Version ist frei zugänglich unter:

\begin{itemize}
	\item \href{https://mathandphysics.de}{\texttt{https://mathandphysics.de}}
	\item \href{https://zenodo.org/communities/christian-copilot}{\texttt{https://zenodo.org/communities/christian-copilot}}
\end{itemize}

Die gedruckte Ausgabe wurde für komfortables Lesen, langfristige Archivierung und den Aufbau einer Open-Access-Bibliothek in Printform erstellt. Mit dem Kauf dieses Buches unterstützen Sie die freie wissenschaftliche Publikation und den Gedanken einer offenen, überprüfbaren Wissenschaft.